\documentclass[10pt]{beamer}

\documentclass[svgnames,smaller]{beamer}

\usetheme{metropolis}

\usepackage[utf8]{inputenc}
\usepackage[russian,french]{babel}

\usepackage[T1]{fontenc}
\usepackage{geometry}
\geometry{
    left=2.5mm,
    right=2.5mm,
    top=0mm,
    bottom=0mm
}
\usepackage{amssymb} % pour \blackdiamond
\usepackage{tikz}
\usepackage{graphicx,animate}

\newcommand{\replacementchar}{%
  \tikz[baseline=-0.6ex]{
    %\node[inner sep=0pt] (d) {\Large$\blackdiamond$};
    \node[inner sep=0pt] (d) {\Large$\blacksquare$};
    \node[white, font=\scriptsize\bfseries] at (d.center) {?};
  }%
}
\usepackage{tabularx}
%\usepackage[style=alphabetic,citestyle=authoryear]{biblatex}
%\addbibresource{bibliographie.bib}

\usepackage{xspace}
\newcommand{\themename}{\textbf{\textsc{metropolis}}\xspace}
\title{L2S4: Édition numérique\\
TD 4: Modéliser et automatiser l'écriture\\
\small
De la combinatoire aux LLM et IAG
}

\date{}
\author{Marine Tiger\\  \quad {marine.tiger@sorbonne-universite.fr\\}}

\institute{Littératures françaises et comparée ED19, CELLF, Sorbonne Université}

\begin{document}
\maketitle
\setbeamertemplate{section in toc}[sections numbered]

\section{Introduction}

\begin{frame}{Problématique esthétique: statut de l'auteur et imaginaire de la création}
    \begin{itemize}
    \item La création renvoie à l'inspiration, à la sacralisation de l'artiste.
    \item L'auteur = un être "inspiré", de "génie" ou qui maîtrise une technique.
\end{itemize}
\centering

\includegraphics[width=0.50\textwidth]{04_Axe1_IAG/Images/prompt.png}\hspace{10mm}

\includegraphics[width=0.40\textwidth]{04_Axe1_IAG/Images/SiROmeNavaitPasChute.png}

\end{frame}

\begin{frame}{Problématique esthétique: statut de l'auteur et imaginaire de la création}
    \begin{block}{}
        \textbf{Principe de "modélisation":} Décomposition d'un objectif en plusieurs étapes, scénarios possibles, afin de produire un programme fonctionnel.\\
        Exemple des applications de rencontre.
    \end{block}
    \begin{block}{Cela implique:}
        \begin{itemize}
            \item Un travail conceptuel majeur \textrightarrow{} comprendre les phénomènes pour pouvoir les décrire et décomposer étape par étape;
            \item Un travail forcément biaisé \textrightarrow{} tenter d'objectivité un phénomène.
        \end{itemize}
    \end{block}

\begin{block}{}
Système probabiliste = réponse basée sur la probabilité de la réponse.
   
\end{block}
\end{frame}

\begin{frame}{La problématique sociale : automatiser/remplacer les métiers (auteurs, éditeurs)}
\begin{columns}[T,onlytextwidth]
\column{0.50\linewidth}
\small
Dans le secteur éditorial, des tâches spécifiques sont aujourd'hui menacées :
\begin{itemize}
    \item la traduction
    \item la révision ortho-typo
    \item la mise en page
\end{itemize}

\column{0.50\linewidth}
\includegraphics[width=0.80\textwidth]{04_Axe1_IAG/Images/IA-traducteurs.png}
\end{columns}
\end{frame}

\begin{frame}{Modélisation et automatisation : deux caractéristiques de l'édition}
\begin{block}{Pour rappel, l'histoire de l'édition c'est:}
    \begin{itemize}
    \item La formalisation des normes et des standard éditoriaux ;
    \item Les techniques d'automatisation
\end{itemize}
\textrightarrow{} Structurent les métiers de l'édition et l'ensemble de la chaîne de production des ouvrages.
\end{block}

\begin{block}{Une modélisation de l'édition = chaîne éditoriale}
    \begin{itemize}
        \item Quelles sont les étapes pour produire un livre ?
        \item À quoi un livre doit-il ressemble (règles de présentation, de mise en forme) ?
    \end{itemize}
    Exemple: la classification Dewey/CDU.
\end{block}
\end{frame}

\begin{frame}{La problématique socio-technique}
\begin{itemize}
    \item Coût environnemental
    \item Coût social, sociétal
    \item Coût politique, éthique
\end{itemize}
\end{frame}

\section{Un débat de société face à un changement paradigmatique majeur}
\begin{frame}{Le péril IA ? La bataille des pionniers}
\begin{block}{}
    \textbf{\textrightarrow{}} Pétition \textit{Pause Giant AI Experiments: An Open Letter} en janvier 2023. Parmi les signataires: Elon Musk, Yoshua Bengio ou encore Steve Wauzniak.\\
\small
\url{https://en.wikipedia.org/wiki/Pause_Giant_AI_Experiments:_An_Open_Letter}
\end{block}
\vspace{10mm}

\centering
\includegraphics[width=0.45\textwidth]{04_Axe1_IAG/Images/Bengio.png}
\includegraphics[width=0.45\textwidth]{04_Axe1_IAG/Images/LeCun.png}
\end{frame}

\begin{frame}{Le péril IA ? La bataille des pionniers}
\begin{columns}[T,onlytextwidth]
\column{0.50\linewidth}
\small
Yoshua Bengio, Yann Le Cun et Geoffrey Hinton ont publié ensemble l'article le plus cité dans le domaine des IA 

\column{0.50\linewidth}
\includegraphics[width=0.80\textwidth]{04_Axe1_IAG/Images/articleIABengio-LeCun.png}
\end{columns}
\end{frame}

\begin{frame}{Le péril IA ? La bataille des pionniers}
\begin{block}{}
    Dans l'article du monde, Yoshua Bengio écrit:\\
    \vspace{3mm}
    \centering
    {\small \textit{"Les résultats restent imprévisibles, et nous ne comprenons pas parfaitement le fonctionnement de ce processus très complexe. Cela pose des questions de sécurité, de transparence, de fiabilité. On s’aperçoit très vite que ces systèmes inventent parfois des réponses, qu’ils peuvent aller dans une direction que n’ont pas choisie les programmeurs. La seule façon de les réorienter, c’est de les récompenser lorsque la réponse est pertinente. Depuis des mois, les concepteurs de ChatGPT travaillent à parer les coups, mais ce n’est pas suffisant."}}\\
\end{block}
\vspace{5mm}
\textrightarrow{} Manque de compréhension sur le fonctionnement de l'IA.\\
\textrightarrow{} Une question \textbf{éthique}.
\end{frame}

\begin{frame}{Le péril IA ? La bataille des pionniers}
\begin{block}{}
    A l'inverse, Yann Le Cun écrit:\\
    \vspace{3mm}
\centering
\small
\textit{"Ces capacités donnent l’impression que le système est intelligent, mais en fait elles restent superficielles. ChatGPT n’est pas intelligent comme peut l’être un humain. C’est un outil de prédiction, qui associe entre eux des mots apparaissant de façon la plus probable dans le corpus qui a servi à l’entraîner, afin de continuer un texte. [...] Ce sont des outils très utiles pour accélérer l’écriture de textes, en particulier de code informatique. Ils peuvent améliorer la productivité de nombreux professionnels, comme les programmeurs, ou les médecins qui veulent rédiger plus vite leurs comptes rendus de consultations. Il faut seulement informer les utilisateurs des limites de l’application : on ne peut pas s’en servir comme moteurs de recherche, ni se reposer sur les informations qu’elle produit sans les vérifier."}
\end{block}
\vspace{5mm}
\textrightarrow{} Selon lui, l'IA peut "conduire à une renaissance de l'humanité, un nouveau siècle des Lumières".
\end{frame}

\begin{frame}{Le péril IA ? La bataille des pionniers}
    \centering
    \includegraphics[width=7cm]{04_Axe1_IAG/Images/The_Cake_is_a_Lie.jpg}
\end{frame}
\begin{frame}{Au fait, c'est quoi les IAG?}
\begin{block}{}
    \textbf{IAG (Intelligence artificielle générative)=} désignent des technologies de création de contenu, de génération de contenus (texte, images, sons), aujourd'hui implémentées dans des outils grand public : copilot, chatGPT, etc.
\end{block}
\vspace{3mm}
\centering
\includegraphics[width=0.65\textwidth]{04_Axe1_IAG/Images/googleNgramIA.png}
\end{frame}

\begin{frame}{"Intelligence artificielle"
une expression problématique et polémique}

\centering
\textbf{Qu'est-ce que l'intelligence?}

\end{frame}

\begin{frame}{Générative ?
}
\begin{block}{}
    \begin{itemize}
        \item Automatiser la production du texte
\item Le mythe du texte qui s'écrit tout seul
\item Une quête formaliste, poétique, ludique...\\
...qui n'a pas attendu ChatGPT !
    \end{itemize}
\end{block}

Rappel: les générateurs de texte=programme qui crée du texte à partir d’un ensemble de règles qui constituent une grammaire et d’un ensemble d’éléments préconstruits qui forment un dictionnaire. 

\end{frame}

\begin{frame}{Les générateurs combinatoires et automatiques}
Repose initialement sur un principe \textbf{combinatoire}.\\
\vspace{5mm}
\textbf{Un générateur combinatoire} = générateur de texte qui combine selon des règles algorithmiques spécifiques des fragments de textes préconstruits. Le générateur va ainsi tenter d’épuiser toutes les possibilités d’une structure à partir de différentes combinaisons des fragments entre eux.
\end{frame}

\begin{frame}{Atelier : étudier la modélisation de l'échange épistolaire amoureux chez Christopher Stratchey}
\centering

\includegraphics[width=0.40\textwidth]{04_Axe1_IAG/Images/atelier-love-letters.png}
\end{frame}

\begin{frame}{Oeuvres génératives}
    \centering
    {\small
    \textit{Ce qui m’intéresse dans la génération, ce n’est pas le texte qui s’affiche. Ce texte-là est un moment comme un autre, on s’en fout. […] Ce qui m’intéresse, c’est cette capacité à produire à l’infini et à générer un univers que je ne suis pas capable de faire. C’est donc un autre substitut qui transmet une pensée qui dit. Peut-être est-ce un fantasme d’éternité.}}\\
\textbf{Jean-Pierre Balpe
\textit{Oeuvre dans un paradigme "contemporain"} (emprunté à l'art contemporain)}

\end{frame}

\section{Les intelligences artificielles : 70 ans d'histoire}
\begin{frame}{Déjà trois printemps et deux hivers
(bientôt trois ?)}

\centering
\includegraphics[width=9cm]{04_Axe1_IAG/Images/frise_ia_genially.png}
\end{frame}

\begin{frame}{Première vague : les années 1950 et les développements de l'informatique}
\textbf{Alan Turing}\\ 
\vspace{5mm}
\begin{columns}[T,onlytextwidth]
\column{0.50\linewidth}
\small
\textrightarrow{} Père de l'informatique.\\
\textrightarrow{} Article fondateur publié en 1936 qui pose les bases du binaire.\\
Rappel : Turing n'a quasiment joué aucun rôle dans la construction effective des ordi.\\ 
\textbf{Machine de Turing =} un automate imaginaire muni d’un programme et pouvant lire et écrire des caractères sur un ruban de longueur illimitée.

\column{0.60\linewidth}
\includegraphics[width=0.80\textwidth]{04_Axe1_IAG/Images/turing.jpg}
\end{columns}

\end{frame}

\begin{frame}{Première vague : les années 1950 et les développements de
l’informatique}

\textbf{Le test de Turing (1950) : une définition de l'intelligence au miroir de la machine}

\begin{columns}[T,onlytextwidth]
\column{0.40\linewidth}
\small
\begin{itemize}
    \item "\textit{Computing Machinery and Intelligence}" (Mind, 1950)
    \item Si on intéragit avec une machine 5 minutes sans réaliser qu'on est face à une machine, alors cette machine peut être qualifiée d'"intelligente"
\end{itemize}
\column{0.50\linewidth}
\includegraphics[width=0.50\textwidth]{04_Axe1_IAG/Images/Turing_intelligence.png}
\end{columns}
\end{frame}

\begin{frame}{Première vague : les années 1950 et les développements de
l’informatique}
\textbf{La conférence de Dartmouth et la proposition de John McCarthy (1956)}\\
\vspace{10mm}
\begin{columns}[T,onlytextwidth]
\column{0.40\linewidth}
\small
= Atelier scientifique organisé durant l'été 1956\\
= Considéré comme l'acte de naissance de l'intelligence artificielle en tant que domaine de recherche autonome\\
\textrightarrow{} Invention du terme \textit{Intelligence artificielle} par John McCarthy.
\column{0.50\linewidth}
\includegraphics[width=0.90\textwidth]{04_Axe1_IAG/Images/darthmouth.jpg}
\end{columns}

\end{frame}

\begin{frame}{Seconde vague : les années 1980 et le développement des "systèmes experts"}

\textbf{Système expert }= programme informatique conçu pour simuler le savoir-faire d’un spécialiste, dans un domaine précis et bien délimité, grâce à l’exploitation d’un certain nombre de connaissances fournies explicitement par des experts du domaine.\\
\vspace{10mm}
\textbf{Exemple de système expert les plus médiatisés} : Garry Kasparov vs Deep Blue.\\
\url{https://www.youtube.com/watch?v=ZIcZymAzifM}

\end{frame}

\begin{frame}{Approche symbolique}
\centering
Les premières conceptions de l'IA s'emploient à programmer la machine pour qu'elle puisse reproduire les formes symboliques et logiques du raisonnement dit naturel. La notion d'"intelligence" est alors elle-même située : il s'agit de transférer à la machine une capacité à raisonner calquée sur le modèle humain.

Mais savons-nous vraiment comment nous pensons ?
    
\end{frame}

\begin{frame}{Seconde vague : les années 1980 et le développement des "systèmes experts"}
\centering
\includegraphics[width=10cm]{04_Axe1_IAG/Images/ModeleExpertSchema.png}
    
\end{frame}

\begin{frame}{Seconde vague : les années 1980 et le développement des "systèmes experts"}

\url{https://www.youtube.com/watch?v=jpCa9D2KPXM}

\begin{block}{}
    \centering
    {\small
    \textit{La réaction de Garry Kasparov immédiatement après ce match est intéressante. Il accuse l’équipe de Deep Blue d’avoir triché, d’avoir fait appel à un joueur humain pour le battre. Pour le champion du monde, aucune machine n’aurait pu à la fois jouer l’excellent 36e coup et commettre cette erreur grossière au 44e. C’est donc forcément un joueur humain qui a dicté ce coup à Deep Blue. Comme beaucoup de non-spécialistes, Kasparov ne croit pas que la machine puisse être à ce point faillible. Il est à mille lieues de la vérité.}}\\
    \textbf{Nicolas Sabouret, \textit{"Pourquoi l’intelligence artificielle se trompe tout le temps"}, The Conversation, 2020.}
\end{block}
\end{frame}

\begin{frame}{La théorie de la complexité et l'approche heuristique}
\textbf{Théorie de la complexité=} étudie formellement le temps de calcul, l'espace mémoire (et plus marginalement la taille d'un circuit, le nombre de processeurs, l'énergie consommée…) requis par un algorithme pour résoudre un problème algorithmique.\\
\vspace{3mm}
\textbf{Exemple des échecs:}
\begin{itemize}
    \item Regarder toutes les parties possibles.
    \item Le mathématicien Claude Shannon a estimé qu’il y a environ 10 puissance 120 parties d’échecs possibles.
    \item Même avec des machines avec des puissances de calcul qui doublerait tous les 2 ans (loi de Moore), impossible d'énumérer toutes ces parties!
\end{itemize}

En résumé: calculer la probabilité de tous les coups possibles est trop complexe.
\end{frame}

\begin{frame}{La théorie de la complexité et l'approche heuristique}
    \centering
    {\small
    \textit{
    Le principe des programmes d’IA est donc de calculer des solutions pas trop mauvaises à des problèmes dont on sait, mathématiquement, qu’ils ne peuvent pas être résolus de façon exacte dans un temps de calcul raisonnable. Et pour cela, il faut accepter de faire parfois des erreurs. De ne pas avoir toujours la meilleure réponse, ou une réponse complètement correcte. C’est pourquoi tout programme d’IA fait forcément des erreurs. C’est inévitable et c’est même ce qui les caractérise.}
    }\\
\textbf{Nicolas Sabouret, \textit{"Pourquoi l’intelligence artificielle se trompe tout le temps"}, The Conversation, 2020}
\end{frame}

\begin{frame}{Limite des systèmes experts, ou le second hiver de l'IA}
\begin{itemize}
    \item Coût de développement trop élevées
    \item Cantonnés à certains domaines industriels: système de signalisation des trains, système de guidage des avions...
\end{itemize}

\centering
{\small
\textit{Qu’est-ce qui n’allait pas dans l’idée d’une machine raisonnant logiquement ? Tout simplement, que le fonctionnement de la pensée humaine est impossible à reproduire. Nous prenons très rarement des décisions à partir de règles de raisonnement que nous saurions expliciter. Nos jugements sont aussi faits d’émotions, d’éléments irrationnels, de spécifications liées au contexte et de toute une série de facteurs implicites ; bref, la décision ne se laisse pas capturer par des règles formalisables.}
}\\
\textbf{Dominique Cardon, Culture numérique, Presses de Science Po, 2019}

\end{frame}

\begin{frame}{Troisième vague : le modèle connexionniste des années 2020}

\begin{block}{}
 \centering
\includegraphics[width=0.90\textwidth]{04_Axe1_IAG/Images/ia-symboliqueVSconnexionniste.jpg}   
\end{block}

Luc Julia dans son ouvrage \textit{L’intelligence artificielle n’existe pas} paru en 2019 aux Éditions First
\begin{itemize}
    \item Ce n’est pas l’intelligence qui caractérise les systèmes d’IA d’aujourd’hui, mais leur capacité de \textbf{reconnaissance grâce à l’apprentissage machine}.
    \item Approche dite "connexionniste", inspirée de \textbf{la structure et du fonctionnement des réseaux de neurones biologiques}.
\end{itemize}

\end{frame}

\begin{frame}{Un concept: les réseaux de neurones}
\begin{columns}[T,onlytextwidth]
\column{0.40\linewidth}
\small
\begin{itemize}
        \item \textbf{\textit{Machine learning}}: Apprentissage à partir de réseaux de neurones artificiels, tente de reproduire le raisonnement humain.
        \item \textbf{\textit{Deep learning}}: Apprentissage à partir de réseaux de neurone profonds pour identifier des structures dans des volumes considérables de données.
    \end{itemize}
\column{0.50\linewidth}
\includegraphics[width=0.80\textwidth]{04_Axe1_IAG/Images/schema_machinelearning_deeplearning.png}
\end{columns}

\end{frame}
\begin{frame}{Un concept: les réseaux de neurones}
\centering
\includegraphics[width=0.60\textwidth]{04_Axe1_IAG/Images/nauronal_network.png}\\
\begin{block}{}
    \textbf{Le modèle connexionniste provoque un changement de paradigme, en nous faisant passer d'une approche modélisante à une approche statistique et probabiliste (fondées sur des modèles malgré tout : les LLM)}
\end{block}
\end{frame}

\begin{frame}{La vectorisation : vers une modélisation spatiale des termes (vs sens)}
    \textbf{Vectorisation=} Convertir des données textuelles (mots, phrases ou documents) en vecteurs numériques.\\

\begin{block}{}
    \centering
\includegraphics[width=0.70\textwidth]{04_Axe1_IAG/Images/wordvectors.jpeg}
\end{block}


\textrightarrow{}Plus les mots sont proches, plus ils ont des chances d'avoir des significations similaires, ou de s'inscrire dans un même contexte sémantique (une même relation contextuelle).

\end{frame}

\begin{frame}{La vectorisation : vers une modélisation spatiale des termes (vs sens)}
\centering
    En conséquence, lorsque vous utilisez une IAG, vous lui posez une question ou vous lui donnez une instruction en vous inscrivant dans un régime de signification. La réponse de l'IAG, en revanche, s'appuiera sur des probabilités, soit sur un calcul de proximité entre des objets vectorisés.
\end{frame}

\begin{frame}{Pour résumer}
    \textbf{Comment aimez-vous apprendre ?}\\
    \begin{itemize}
        \item \textbf{Approche symbolique} : vous préférez que l'on vous donne des consignes précises, après avoir suivi un cours magistral en amont vous présentant des règles de composition d'une dissertation, que vous allez ensuite vous-même appliquer.
        \item \textbf{Approche connexionniste} : vous préférez lire plusieurs exemples de dissertations, en tâchant de comprendre vous-mêmes comment elles fonctionnent pour les imiter.
    \end{itemize}

\end{frame}

\section{La problématique des outils fondés sur les IAG}
\begin{frame}{ChatGPT, l'IA générative de texte d'OpenAI}
\begin{block}{}
    \textbf{ChatGPT ?
    GPT (GPT3, GPT4, etc.)}\\

\begin{itemize}
    \item \textit{Generative Pre-trained Transformer}= modèle géant de prédiction de texte entraîné par OpenAI sur 500 milliards de mots (pour GPT3).
    \item Modèle encyclopédique
    \item Appartient à la famille des LLM= entraîné sur d'immenses quantités de données.
\end{itemize}
\end{block}
    
\begin{block}{}
\textbf{Chat (InstructGPT)}
    \begin{itemize}
        \item Une interface conversationnelle
        \item Fonctionne à partir des requêtes de l'usager : "le prompt"
    \end{itemize}
\end{block}
\end{frame}

\begin{frame}{Avec qui je "converse", lorsque j'échange avec ChatGPT ?
}
\textbf{1. Un agent conversationnel intégré à une interface au design conçu pour masquer la technique}
\vspace{5mm}
\begin{columns}[T,onlytextwidth]
\column{0.40\linewidth}
\small
\begin{itemize}
        \item le tout premier chatbot psychothérapeute (1966) conçu par Joseph Weizenbaum. Version Javascript par Michael Wallace et George Dunlop
    \end{itemize}
\column{0.50\linewidth}
\includegraphics[width=0.80\textwidth]{04_Axe1_IAG/Images/Eliza.png}
\end{columns}
    
\end{frame}

\begin{frame}{Avec qui je "converse", lorsque j'échange avec ChatGPT ?}
    \textbf{2. Un LLM : un corpus de texte construit à partir du big data}

    \vspace{5mm}

\includegraphics[width=0.80\textwidth]{04_Axe1_IAG/Images/2021-Alan-D-Thompson-GPT-3-datasets-by-effective-size-v2.png}
\end{frame}

\begin{frame}{Problématique des origines des données d'entraînement}
OpenAI ne publie pas la totalité des spécifications sur la construction de ses modèles.\\

Voici les trois catégories principales de données selon OpenAI :
\begin{itemize}
    \item \textbf{Données accessibles publiquement sur Internet}\\
    \textrightarrow{} Quid des données personnelles?
    \item \textbf{Données sous licence ou partenaires tiers}\\
    \textrightarrow{} Contenu auquel OpenAI a légalement accès
    \item Données fournies ou générées avec participation humaine\\
    \textrightarrow{} Données annotées, modèles entraînés, fine-tuning. digital labour.
\end{itemize}

Question de la neutralité des llm.
\end{frame}

\begin{frame}{Les biais de l'IA et des algorithmes}
    \textbf{Biais de l'IA =} Désignent les systèmes d’IA qui produisent des résultats biaisés répétant les préjugés humains au sein d’une société\\

\includegraphics[width=0.60\textwidth]{04_Axe1_IAG/Images/UNESCO.png}

\end{frame}

\begin{frame}{ChatGPT et les biais de genre}

Isabelle Collet, professeure en sciences de l’éducation à l’Université de Genève et spécialiste des questions de genre dans le numérique, dans son livre \textit{Le numérique est l’affaire de toutes} distingue trois principaux risques :
\begin{itemize}
    \item \textbf{L’apprentissage des stéréotypes}. Ex de Wikipédia.
    \item \textbf{La discrimination dans l’aide à la décision}. Renforcement des inégalités déjà présentes dans la société, par exemple en matière d’embauche, de rémunération, de promotion, d’octroi de prêts ou encore d’orientation scolaire.
    \item \textbf{Le manque de diversité dans les métiers du numérique.} Invisibilisation d'une partie des besoins et des expériences de la population.
\end{itemize}
\end{frame}

\begin{frame}{Midjourney et les stéréotypes racistes}
\begin{columns}[T,onlytextwidth]
\column{0.50\linewidth}

\includegraphics[width=0.80\textwidth]{04_Axe1_IAG/Images/franceinter.png}\\
\small
\href{https://www.radiofrance.fr/franceinter/tags-dechets-batiments-delabres-quand-l-ia-reproduit-les-stereotypes-sur-les-banlieues-3334839}{Podcast sur France Inter}.
\column{0.50\linewidth}

\includegraphics[width=0.90\textwidth]{04_Axe1_IAG/Images/refecltions_before_storm.png}\\
\small
Alenichev, Arsenii et al., \href{https://www.thelancet.com/journals/langlo/article/PIIS2214-109X(23)00329-7/fulltext?utm_source=generationia.flint.media&utm_medium=referral&utm_campaign=dall-e-midjourney-les-stereotypes-raciaux-sont-encore-tenaces-dans-l-ia}{\textit{Reflections before the storm: the AI reproduction of biased imagery in global health visuals}}, 2023

\end{columns}
\end{frame}

\begin{frame}{Quid du coût écologique?}

Ce qu'il faut pour entraîner des modèles d'IAG:
\begin{itemize}
\item Construire des serveurs
    \item De l'eau pour refroidir les serveurs
    \item De l'électricité
\end{itemize}
\begin{block}
\centering
\includegraphics[width=0.60\textwidth]{04_Axe1_IAG/Images/mistral_ecolo.png}    
\end{block}


Pour l'entraînement de Mistral Large 2, après 18 mois d'utilisation:
\begin{itemize}
    \item 20,4 ktCO₂e
\item 281 000 m3 d'eau consommée
\item 660 kg Sb eq (unité standard pour l'épuisement des ressources).
\end{itemize}

\end{frame}

\begin{frame}{Pourquoi ChatGPT n'est PAS un moteur de recherche ?}
    
ChatGPT n’a aucune compréhension de votre question ni même de sa réponse, qui résulte d'un calcul statistique croisant le contexte de votre requête et les différentes données qu’il a à sa disposition.
\end{frame}

\begin{frame}{Attention aux hallucinations!}
    \textbf{Exemple: invention d'une source bibliographique.}\\
    \vspace{5mm}
    \centering
    \includegraphics[width=0.70\textwidth]{04_Axe1_IAG/Images/gptHalluBoullier.png}
\end{frame}

\begin{frame}{ChatGPT search, un moteur de réponse}
\begin{itemize}
    \item Un service qui nous propose une synthèse de certaines sources en fonction de notre requête
    \item Une liste de sources (critères de sélection non précisés)
    \item Pas de rémunération des sources
    \item Une logique de plus en plus prescriptive
    \item Quid de la compétence d'évaluation et d'esprit critique ?
\end{itemize}
    
\end{frame}

\begin{frame}{Les IAG dans le secteur des industries culturelles et créatives : une communauté divisée}
\begin{itemize}
    \item Les modèles entraînés respectent-ils le droit d'auteur?
    \item L'IA est-elle un outil de création comme un autre, ou sert-elle une imposture ?
    \item En améliorant la productivité, dans quelle mesure les IAG menacent-elle les professionnels ?
    \item Le recours à l'IAG vaut-il vraiment le coup (engagement : enjeux éthiques et écologiques) ?
\end{itemize}

\end{frame}

\section{Les industries éditoriales et littéraires}
\begin{frame}{Les industries éditoriales et
littéraires}
 Une IA qui va intégrer par contre des outils d'édition :
\begin{itemize}
    \item révision de texte
    \item gestion des plannings
    \item communication
    \item Globalement, une industrie qui apprécie une forme d'artisanat et qui se montre assez réfractaire.
\end{itemize}


Cas particulier : illustrateurs.
\end{frame}

\begin{frame}{Cas n°1 : Mathis et la forêt des possibles}
\begin{block}{}
    \centering
    \includegraphics[width=0.60\textwidth]{04_Axe1_IAG/Images/mathis-Foret.png}
\end{block}
\begin{itemize}
    \item Collaboration entre les PUR et la maison d'édition Locus Solus
    \item Jiri Benovsky est spécialiste en métaphysique, en philosophie de l’esprit et en esthétique. Pas illustrateur.
    \item Son premier album jeunesse, réalisé par IA.
\end{itemize}

\end{frame}

\begin{frame}{Un bon prompteur peut-il remplacer un illustrateur ?}
\begin{block}{}
\textrightarrow{} Qu'est-ce que c'est que le geste de création ?
    \textrightarrow{} Vieux débat sur les conditions et les compétences pour mériter le statut d'artiste.
\end{block}
\centering
    \includegraphics[width=0.60\textwidth]{04_Axe1_IAG/Images/mathisEtLaForetDesPossibles.png}
\end{frame}

\begin{frame}{La problématique des données et des modèles entraînés}
On a toujours utilisé des outils pour nous aider dans notre travais mais est-ce qu'une IA est un outil comme un autre ? \\
Exemple: des IA entraînées sur des oeuvres protégées par le droit d'auteur.\\
\vspace{3mm}
\begin{block}{}
    \centering
    \includegraphics[width=0.70\textwidth]{04_Axe1_IAG/Images/author_suing_gpt.png}
\end{block}

\textrightarrow{} Crainte des créateur.rices de l'accaparement de la propriété intellectuelle, baffouée par les modèles.
\end{frame}

\begin{frame}{Constituer et entraîner un VML : un geste technique, un geste auctorial ?}
\begin{itemize}
    \item Constituer son modèle et l'entraîner : la réalité des pionniers
    \item Le développement vers des outils grand public menace de faire disparaître la technicité
\end{itemize}
Selon Frédéric Maupomé, "\textit{Les IA n’inventent rien : il leur faut plagier la création humaine, en s’inspirant de centaines de millions d’images et métadonnées aspirées automatiquement sur internet.}"

\end{frame}

\begin{frame}{Un droit à l'\textit{opt-out} peu efficace}

\textbf{\textit{opt-out}}= désigne la possibilité pour une personne, une organisation ou un créateur de refuser que ses données, œuvres ou contenus soient utilisés pour l’entraînement, l’évaluation ou le fonctionnement d’un modèle d’IA.\\
\vspace{5mm}

\centering
{\small

\textit{Un calvaire ! C'est image par image, il faut prouver qu'on est l'ayant droit, recommencer parfois à chaque copie présente dans la base, (ex. Have i been trained). Et recommencer pour chaque base... Sans preuve d'effet (et l'opt-out n'est pas rétroactif), si on est dans la base,c'est trop tard. Les restrictions mises en place par les entreprise d'IA (par exemple on ne peut pas demander d'image dans le style de Steve mc curry à stable diffusion, il y a un filtre sur son nom) sont facilement contournées (par exemple en écrivant mal les noms).}}\\
\textbf{ Positionnement de la Ligue des auteurs professionnels}
\end{frame}

\begin{frame}{Atelier: Du programme au prompt : comprendre les stratégies de génération de récits à l’heure des LLM}
\centering
    
\includegraphics[width=4cm]{04_Axe1_IAG/Images/atelier-prompt.png}
\end{frame}

\begin{frame}{Au-delà des industries culturelles et créatives...}
\begin{itemize}
    \item Des enjeux éthiques : entraînement des données et \textit{digital labour}
    \item Des enjeux politiques : nouveaux monopoles, vectofascisme
    \item Des enjeux écologiques : coût d'entraînement des modèles (calcul, stockage)
\end{itemize}
    
\end{frame}

\begin{frame}{L'IA Act: un premier pas en faveur de la régulation}
   \begin{columns}[T,onlytextwidth]
\column{0.50\linewidth}
\small
\textit{Puisqu’il sera impossible de garantir l’absence de discrimination sous tous les angles, des choix seront nécessaires, et des seuils de tolérance devront être définis. Se pose alors la question : quel est le niveau acceptable d’erreurs et de quel type d’erreurs parle-t-on ? Ce qui revient à se demander : quel est le niveau « acceptable » de risque de discrimination ?}\\ 
\textbf{Mélanie Gornet et Winston Maxwell, "\textit{Intelligence artificielle: normes techniques et droits fondamentaux, un mélange risqué}", The Conversation, 2022.}
\column{0.50\linewidth}
\centering
\includegraphics[width=0.90\textwidth]{04_Axe1_IAG/Images/AI-ACT.png}
\end{columns}
\end{frame}


\section{Conclusion}
\begin{frame}{Conclusion}
\centering
Au-delà des effets de fascination - largement influencés par notre imaginaire fictionnel de la machine - nous commençons tout juste à prendre la mesure des effets de l'IA sur la création, la créativité, et les conditions de travail dans les industries culturelles. Si tout le monde s'accorde sur le besoin de réguler, les leviers à (dés)activer font largement débat. Assistant ou remplaçant, l'outil d'IA doit en tous les cas être bien compris et transparent afin que d'en faire un usage bien conscient. Le véritable coût économique, social et écologique de la dernière génération d'IA reste encore inconnu, et l'on peut justement s'interroger sur les développements futur de cette technologie : un autre hiver de l'IA est-il à prévoir ? Ou sommes-nous définitivement entrés dans un changement de paradigme ?
    
\end{frame}
\end{document}